\documentclass[11pt]{article}
\usepackage[a4paper,margin=1.9cm]{geometry}
\usepackage[T1]{fontenc}
\usepackage{textcomp}
\usepackage[spanish]{babel}
\usepackage{amsmath,amssymb}
\usepackage{enumitem}
\usepackage{tikz}
\usetikzlibrary{calc,arrows.meta,patterns,decorations.pathmorphing}
\usepackage{xcolor}
\usepackage{multicol}
\usepackage{parskip}
\usepackage{lmodern}
\renewcommand{\familydefault}{\sfdefault}

\definecolor{unalblue}{RGB}{31,73,124}
\definecolor{unalblue2}{RGB}{70,120,180}

\newlist{opciones}{enumerate}{1}
\setlist[opciones]{label=\textbf{(\alph*)},itemsep=1pt,topsep=2pt,leftmargin=1.8em,font=\normalfont}
\setlist[enumerate,1]{label=\textbf{\arabic*.},itemsep=6pt,topsep=4pt,leftmargin=1.6em,font=\normalfont}
\setlist[enumerate,2]{label=\textbf{\alph*)},itemsep=4pt,topsep=3pt,leftmargin=1.6em,font=\normalfont}
\setlist[enumerate,3]{label=\textbf{\roman*)},itemsep=2pt,topsep=2pt,leftmargin=1.4em,font=\normalfont}

\newcommand{\vect}[2]{\begin{pmatrix}#1\\#2\end{pmatrix}}
\newcommand{\vectiii}[3]{\begin{pmatrix}#1\\#2\\#3\end{pmatrix}}

\newcommand{\examheader}{%
\noindent\begin{tikzpicture}
\node[fill=unalblue, text width=\dimexpr\linewidth-16pt\relax, inner sep=8pt, rounded corners=3pt, align=left, text=white] (hdr) {%
  \begin{minipage}[c]{0.66\linewidth}
    {\Large\bfseries \examsubject}\\[2pt]
    {\small \examtype\ \textbullet\ \textcolor{white!85!unalblue}{\examsubtitle}}
  \end{minipage}\hfill
  \begin{minipage}[c]{0.28\linewidth}
    \raggedleft{\scriptsize Escuela de Matem\'aticas\\[-1pt] Universidad Nacional\\[-1pt] de Colombia, Sede Medell\'in}
  \end{minipage}%
};
\end{tikzpicture}\\[7pt]
}
\newcommand{\examfooter}{%
  \vfill
  \rule{\linewidth}{0.4pt}\\[1pt]
  {\scriptsize\textit{Transcripci\'on digital --- MathUNAL}\hfill\textcolor{unalblue}{\textbf{\thepage}}}
}
\pagestyle{empty}

\newcommand{\examsubject}{ECUACIONES DIFERENCIALES}
\newcommand{\examtype}{Tercer Parcial}
\newcommand{\examsubtitle}{Semestre 01-2018, 31 de Mayo de 2018}

\begin{document}

\examheader

\vspace{4pt}
\noindent
\begin{minipage}[t]{0.55\linewidth}
{\bfseries\large Tercer Examen Parcial Ecuaciones Diferenciales}\\[2pt]
Puntaje. S\'olo para uso Oficial\\[4pt]
\begin{tabular}{|c|c|c|c|c|c|}
\hline
1 (/30) & 2 (/15) & 3 (/30) & 4 (/25) & Total & Nota\\
\hline
 & & & & & \\
\hline
\end{tabular}
\end{minipage}%
\hfill
\begin{minipage}[t]{0.4\linewidth}
Nombre: \dotfill\\[6pt]
C\'edula: \dotfill\\[4pt]
Sal\'on: \dotfill \hfill Grupo: TODOS\\[10pt]
{\large Examen N\'umero: \dotfill}\\[10pt]
Firma: \dotfill
\end{minipage}

\examfooter
\newpage

\examheader

\begin{multicols}{2}
\begin{enumerate}
\item \textbf{(30\%)} Resuelva el siguiente problema de valor inicial en el intervalo $[0,\infty)$:
\[
y'' + y = f(t); \quad y(0)=0, \quad y'(0)=0,
\]
donde
\[
f(t) = \begin{cases} t & \text{si } 0 \le t < 1, \\ 1 & \text{si } 1 \le t. \end{cases}
\]
Nota: en este ejercicio se calificar\'a el procedimiento.

Sugerencia: puede usar que
\[
\frac{1}{s^2(s^2+1)} = \frac{1}{s^2} - \frac{1}{s^2+1}.
\]

\columnbreak

\item \textbf{(15\%)} Considere el sistema no homog\'eneo:
\[
\mathbf{X}'(t) = \mathbf{A}\mathbf{X}(t) + \begin{pmatrix} 1 \\ -1 \end{pmatrix} e^{-t}.
\]
Si se sabe que $\Phi(t) = \begin{pmatrix} e^t & te^t \\ -e^t & \left(\frac{1}{2}-t\right)e^t \end{pmatrix}$ es una matriz fundamental del sistema homog\'eneo asociado, encuentre una soluci\'on particular del sistema no homog\'eneo dado en $(-\infty,\infty)$. Nota: en este ejercicio se calificar\'a el procedimiento.
\end{enumerate}
\end{multicols}

\examfooter
\newpage

\examheader

\begin{enumerate}
\setcounter{enumi}{2}
\item \textbf{(30\%)} En los literales a), b), c) y d), complete seg\'un corresponda. S\'olo se calificar\'a respuesta, no es necesario que escriba el procedimiento.
\begin{enumerate}[label=\alph*)]
\item \textbf{(7\%)} $\displaystyle \mathcal{L}\left\{\int_0^t e^{-2\tau}\tau^4\,d\tau\right\} = $ \dotfill .
\item \textbf{(8\%)} Si $\displaystyle F(s) = \frac{2e^{-\pi s}}{7s+2}$, entonces $\mathcal{L}^{-1}\{F(s)\} = $ \dotfill .
\item \textbf{(8\%)} Si $F(s) = \ln(s-3)$, entonces $\mathcal{L}^{-1}\{F(s)\} = $ \dotfill .
\item \textbf{(7\%)} $\mathcal{L}\{(t\cos t)\,\delta(t-\pi)\} = $ \dotfill .
\end{enumerate}
\end{enumerate}

\examfooter
\newpage

\examheader

\begin{multicols}{2}
\begin{enumerate}
\setcounter{enumi}{3}
\item{}
\begin{enumerate}[label=\alph*)]
\item \textbf{(10\%)} Halle la soluci\'on del P.V.I.:
\[
\mathbf{X}'(t) = \mathbf{A}\mathbf{X}(t), \quad \mathbf{X}(0) = \begin{pmatrix} 1 \\ 0 \end{pmatrix}
\]
en $(-\infty,\infty)$, sabiendo que $\mathbf{A} \in \mathbb{R}_{2\times 2}$ es una matriz de entradas constantes, $r=-2+i$ es un valor propio de $\mathbf{A}$, y un vector propio correspondiente a $r$ es $\mathbf{K} = \begin{pmatrix} 1-i \\ 1 \end{pmatrix}$.

\columnbreak

\item \textbf{(15\%)} Halle la soluci\'on general del sistema lineal homog\'eneo:
\[
\mathbf{X}'(t) = \begin{pmatrix} -3 & -1 \\ 9 & 3 \end{pmatrix}\mathbf{X}(t)
\]
en el intervalo $(-\infty,\infty)$, si se sabe que $r=0$ es un valor propio repetido (de multiplicidad algebraica 2) de la matriz de coeficientes, y que sus vectores propios correspondientes son todos m\'ultiplos escalares del vector $\mathbf{K} = \begin{pmatrix} 1 \\ -3 \end{pmatrix}$.
\end{enumerate}
\end{enumerate}
\end{multicols}

\examfooter

\end{document}
